\documentclass[11pt]{article}

% --- UNIVERSAL PREAMBLE BLOCK ---
% Set 4:5 Aspect Ratio (Standard LinkedIn Portrait Carousel 1080x1350 equivalent)
\usepackage[paperwidth=8in, paperheight=10in, top=1.2in, bottom=1.2in, left=1in, right=1in]{geometry}
\usepackage{fontspec}
\usepackage[english, bidi=basic, provide=*]{babel}

\babelprovide[import, onchar=ids fonts]{english}

% Set premium typography using Noto Sans system fonts
\babelfont{rm}{Noto Sans}
\babelfont{sf}{Noto Sans}
\babelfont{tt}{Noto Sans Mono}

% --- PACKAGES FOR MODERN DESIGN & MATHS ---
\usepackage{amsmath}
\usepackage{amssymb}
\usepackage{xcolor}
\usepackage{graphicx}
\usepackage{listings}
\usepackage{mdframed}
\usepackage{tikz}
\usetikzlibrary{shapes.geometric, arrows.meta, positioning, shadows.blur}
\usepackage{tcolorbox}
\tcbuselibrary{listings, skins, breakable}
\usepackage{fancyhdr}
\usepackage{enumitem}

% --- DESIGN TOKENS (Stripe & Linear Inspired) ---
\definecolor{brandnavy}{HTML}{0A2540}    % Deep Primary Navy
\definecolor{brandgray}{HTML}{425466}    % Slate Gray
\definecolor{brandlight}{HTML}{F8F9FA}   % Light Soft Gray
\definecolor{brandblue}{HTML}{635BFF}    % Accent Classic Blue
\definecolor{brandgreen}{HTML}{00D4B2}   % Success/Muted Green
\definecolor{brandorange}{HTML}{F25C05}  % Accent Orange
\definecolor{brandred}{HTML}{E5484D}     % Warning Red

% Premium Dark Theme for Code Blocks
\definecolor{codebackground}{HTML}{141419}
\definecolor{codekeyword}{HTML}{FF7B72}
\definecolor{codestring}{HTML}{A5D6FF}
\definecolor{codecomment}{HTML}{8B949E}
\definecolor{codetext}{HTML}{E6EDF3}
\definecolor{codeframe}{HTML}{30363D}

% --- ZERO PARAGRAPH INDENT ---
\setlength{\parindent}{0pt}
\setlength{\parskip}{1.2em}

% Custom TColorBox for styled content blocks
\newtcolorbox{infobox}[1][]{
    colback=brandlight, colframe=brandblue, arc=5pt, boxrule=0.8pt,
    left=15pt, right=15pt, top=12pt, bottom=12pt, #1
}

\newtcolorbox{darkbox}[1][]{
    colback=brandnavy, colframe=brandnavy, arc=5pt, boxrule=0.8pt,
    left=15pt, right=15pt, top=12pt, bottom=12pt, 
    coltext=white, #1
}

% --- HEADER & FOOTER CONFIGURATION ---
\pagestyle{fancy}
\fancyhf{}
\renewcommand{\headrulewidth}{0pt}
\renewcommand{\footrulewidth}{0.5pt}

\fancyhead[L]{\fontsize{9}{11}\selectfont\color{brandgray}\MakeUppercase{First Principles Deep Dive}}
\fancyhead[R]{\fontsize{9}{11}\selectfont\color{brandblue}\textbf{AI ARCHITECTURE}}

\fancyfoot[L]{\fontsize{8}{10}\selectfont\color{brandgray}
Author: Dibayendu Mukherjee \hspace{0.7em}$\cdot$\hspace{0.7em}
\mbox{AI Research Engineer \hspace{0.7em}\textbullet\hspace{0.7em} Systems Engineer}}

\fancyfoot[R]{\fontsize{8}{10}\selectfont\color{brandgray}Slide \thepage}

\begin{document}

% ==========================================================
% SLIDE 1: COVER
% ==========================================================
\thispagestyle{empty}
\vspace*{0.8in}

\begin{center}

    {\fontsize{14}{16}\selectfont\color{brandblue}\textbf{
    THE FULL TRANSFORMER ARCHITECTURE
    }}

    \vspace{0.8in}

    {\fontsize{40}{48}\selectfont\color{brandnavy}\bfseries
    How does AI\\
    understand\\
    true meaning?
    }

    \vspace{0.6in}

    {\fontsize{18}{24}\selectfont\color{brandgray}
    Self-Attention, Multi-Head Spaces,\\
    Cross-Attention, and GPU Tradeoffs.
    }

    \vfill

    \begin{minipage}{0.8\textwidth}
        \centering
        \rule{0.6\linewidth}{0.5pt}

        \vspace{0.2in}

        {\fontsize{12}{14}\selectfont\color{brandnavy}
        \textbf{Dibayendu Mukherjee}
        }

        \vspace{0.05in}

        {\fontsize{10}{12}\selectfont\color{brandgray}
        AI Research Engineer \quad $\cdot$ \quad Systems Engineer
        }

    \end{minipage}

\end{center}

\newpage

% ==========================================================
% SLIDE 2: THE PROBLEM
% ==========================================================
\vspace*{0.4in}

\begin{center}
{\fontsize{26}{30}\selectfont\color{brandnavy}\bfseries 
The Sequential Bottleneck}
\end{center}

\vspace{0.3in}

{\large\color{brandgray}
Before Transformers, NLP relied heavily on recurrent models (RNNs/LSTMs).
}

\vspace{0.15in}

But reading word-by-word introduced two massive flaws:

\textbf{1. Memory Degradation:} By the end of a long sentence, the network often forgot the beginning.

\textbf{2. The Hardware Bottleneck:} Token $N$ had to wait for Token $N-1$ to finish processing. 

\vspace{0.15in}

\begin{darkbox}

{\large\textbf{The Realization}}

\vspace{0.05in}

As long as we processed sequences one step at a time, we could never fully utilize the massive parallel computing power of modern GPUs.

\begin{center}
\textit{We needed a way to process everything at once.}
\end{center}

\end{darkbox}

\vfill
\newpage

% ==========================================================
% SLIDE 3: THE INTUITION
% ==========================================================
\section*{The Intuition: The Cocktail Party}

\vspace{0.15in}

Imagine you are at a crowded cocktail party and you overhear the word \textbf{"apple"}.

Are they talking about the fruit or the tech company?

\vspace{0.2in}

Your brain doesn't read the conversation sequentially to figure it out. It instantly scans the surrounding words for context clues.

\vspace{0.2in}

\begin{minipage}[t]{0.45\textwidth}
    {\large\color{brandblue}\textbf{Context A: Technology}}

    \begin{itemize}[leftmargin=*, itemsep=0.1in, label={\color{brandblue}\textbullet}]
        \item ``Silicon''
        \item ``Macbook''
        \item ``Steve Jobs''
    \end{itemize}

\end{minipage}
\hfill
\begin{minipage}[t]{0.45\textwidth}

    {\large\color{brandorange}\textbf{Context B: Fruit}}

    \begin{itemize}[leftmargin=*, itemsep=0.1in, label={\color{brandorange}\textbullet}]
        \item ``Pie''
        \item ``Orchard''
        \item ``Crisp''
    \end{itemize}

\end{minipage}

\vspace{0.3in}

\begin{infobox}

This is exactly how \textbf{Self-Attention} works. 

Instead of reading left-to-right, every single word looks at every other word in the sentence simultaneously, asking:

\begin{center}
\textbf{"Which of these words modifies my meaning?"}
\end{center}

\end{infobox}

\vfill
\newpage

% ==========================================================
% SLIDE 4: THE SOLUTION (Q, K, V)
% ==========================================================

\section*{The Mechanism: Q, K, and V}

\vspace{0.15in}

To achieve this, the Self-Attention mechanism maps every word into three distinct vectors:

\vspace{0.15in}

\begin{darkbox}[top=8pt, bottom=8pt, left=12pt, right=12pt]
{\large\color{brandlight}\textbf{1. The Query (Q)}}

\vspace{0.05in}
{\fontsize{10}{12}\selectfont
\textit{"What am I looking for?"} \\
This represents the current word that is trying to gather context from the rest of the sentence.
}
\end{darkbox}

\vspace{0.1in}

\begin{darkbox}[top=8pt, bottom=8pt, left=12pt, right=12pt]
{\large\color{brandlight}\textbf{2. The Key (K)}}

\vspace{0.05in}
{\fontsize{10}{12}\selectfont
\textit{"What do I have?"} \\
This represents the information or "tags" that every other word is holding. The network checks if a Key matches the Query.
}
\end{darkbox}

\vspace{0.1in}

\begin{darkbox}[top=8pt, bottom=8pt, left=12pt, right=12pt]
{\large\color{brandlight}\textbf{3. The Value (V)}}

\vspace{0.05in}
{\fontsize{10}{12}\selectfont
\textit{"What is my actual meaning?"} \\
If a Query and a Key match well, the network pulls in this Value to update the original word's context.
}
\end{darkbox}

\vfill
\newpage

% ==========================================================
% SLIDE 5: THE MATH
% ==========================================================
\section*{Breaking Down the Math}

\vspace{0.05in}

{\fontsize{10}{12}\selectfont\color{brandgray}
It looks intimidating, but it is just a remarkably elegant way to compute relationships.
}

\vspace{0.15in}

\begin{center}
{\Large $\text{Attention}(Q, K, V) = \text{softmax}\left(\frac{QK^T}{\sqrt{d_k}}\right)V$}
\end{center}

\vspace{0.15in}

\begin{itemize}[leftmargin=*, itemsep=0.15in, label={\color{brandblue}\textbullet}]
    \item \textbf{$QK^T$ (Dot Product):} We multiply the Queries by the Keys. The dot product measures \textit{similarity}. If word A is highly relevant to word B, this score explodes.
    \item \textbf{$\sqrt{d_k}$ (Scaling):} We divide by the square root of the dimension size to keep the gradients stable during backpropagation.
    \item \textbf{softmax:} This converts the raw scores into percentages (weights that sum to 1). Now we know exactly how much "attention" word A should pay to word B.
    \item \textbf{$\times V$ (The Update):} We multiply those percentages by the actual Values, creating a new, rich, context-aware representation of our word.
    \item \textbf{Spaces of $Q$, $K$, and $V$:} 
         $Q$ and $K$ are projected into the \textit{same vector space}, so their dot product is meaningful — it measures alignment between "what I’m looking for" (Query) and "what I contain" (Key).  
        $V$ lives in the \textit{embedding space}, carrying the semantic content of the words. The attention weights (from $QK^T$) decide how much of each Value contributes to the final representation.
\end{itemize}


\vfill
\newpage

% ==========================================================
% SLIDE 6: MULTI-HEAD ATTENTION
% ==========================================================

\section*{Multi-Head Attention}

\vspace{0.1in}

Why stop at a single perspective? 

Instead of one attention function, Transformers project Q, K, and V into \textbf{multiple lower-dimensional subspaces}.

\vspace{0.15in}

\begin{minipage}[t]{0.47\textwidth}

{\large\color{brandorange}\textbf{The Theory}}

\vspace{0.05in}
We say that multiple heads allow the model to capture different relationships (syntactic vs semantic vs positional) simultaneously.

\end{minipage}
\hfill
\begin{minipage}[t]{0.47\textwidth}

{\large\color{brandblue}\textbf{The Reality}}

\vspace{0.05in}
The model is never explicitly told to split the work! It just discovers these specializations organically through gradient descent.

\end{minipage}

\vspace{0.2in}

\begin{infobox}

\textbf{An Emergent Property}

Separate parameter subspaces naturally increase representational capacity. We didn't hand-engineer syntactic rules into the heads; the network learned that specializing was the most efficient way to minimize the loss function.

\vspace{0.15in}

\textbf{Example:}  
- One head might learn that \textit{“eat”} (verb) strongly attends to \textit{“apple”} (noun), capturing syntactic structure.  
- Another head might focus on semantic roles, linking \textit{“dog”} with \textit{“bark”}.  
- Yet another head could emphasize positional context, such as connecting \textit{“yesterday”} with the verb tense.

These heads are not manually assigned — the model discovers these verb–noun and semantic relationships organically.
\end{infobox}


\vfill
\newpage

% ==========================================================
% SLIDE 7: CROSS-ATTENTION
% ==========================================================

\section*{The Decoder's Magic:\\Cross-Attention}

\vspace{0.1in}

While self-attention looks inward, cross-attention looks outward. It is the bridge between the Encoder and the Decoder. 

\vspace{0.15in}

\begin{darkbox}
\textbf{Self-Attention:} \textit{"What in my own past should I pay attention to?"}

\vspace{0.05in}
\textbf{Cross-Attention:} \textit{"What in the external context should I incorporate?"}
\end{darkbox}

\vspace{0.15in}

In the Decoder, the \textbf{Queries} come from the target sequence (the words we are generating).

But the \textbf{Keys} and \textbf{Values} come directly from the Encoder's final output (the source sentence). 

{\small (Recall: the \textbf{Encoder} compresses the source input into a rich representation, while the \textbf{Decoder} expands that representation into the target output. Source: Previous AE Post.)}

\vspace{0.15in}

Every single token generated dynamically asks the entire source sentence: \textbf{"Which part of you is relevant to what I'm trying to say next?"} This creates a differentiable, dynamic alignment that powers modern translation.

\vspace{0.15in}

\textbf{Example:}  
Suppose the source sentence (Encoder input) is: \textit{"Le chat dort"} (French).  
- The Decoder is trying to generate the English translation.  
- When predicting the word \textit{"cat"}, the Query comes from the partially generated target sequence.  
- The Keys and Values come from the Encoder’s representation of \textit{"Le chat dort"}.  
- Cross-attention aligns the Query with the Encoder’s output, discovering that \textit{"chat"} is the most relevant part.  
- This dynamic alignment ensures the Decoder picks the correct word \textit{"cat"} in context.



\vfill
\newpage

% ==========================================================
% SLIDE 8: THE ENGINEERING TRADEOFF
% ==========================================================

\section*{The Engineering Genius}

\vspace{0.1in}

The brilliance of the "Attention Is All You Need" paper wasn't just linguistic. It was an engineering masterclass.

\vspace{0.15in}

\begin{minipage}[t]{0.47\textwidth}

{\large\color{brandorange}\textbf{The Cost}}

\vspace{0.05in}

Memory complexity became $O(N^2)$. 

Because every word compares itself to every other word, doubling the sentence length quadruples the memory required. 

This is why LLMs struggle with infinite context windows.

\end{minipage}
\hfill
\begin{minipage}[t]{0.47\textwidth}

{\large\color{brandblue}\textbf{The Reward}}

\vspace{0.05in}

Complete Parallelization.

There are no sequential loops. The entire computation is just massive matrix multiplication happening across thousands of GPU cores simultaneously.

\end{minipage}

\vspace{0.2in}

\begin{infobox}

\textbf{The Real Breakthrough}

By shifting the bottleneck from sequential processing to memory bandwidth, transformers unlocked the ability to train on terabytes of data in parallel. 

They traded algorithmic elegance for hardware utilization.

\vspace{0.15in}

\textbf{Example:}  
- Imagine a sentence with 1,000 words.  
- Self-attention requires each word to compare itself with all 999 others, about 1,000,000 comparisons ($N^2$).  
- That’s the \textbf{cost}: memory usage explodes as sequence length grows.  
- But the \textbf{reward}: all those comparisons can be done at once using GPU matrix multiplication. Instead of whispering word-by-word like an RNN, it’s like blasting the whole sentence through thousands of speakers simultaneously.  
\end{infobox}


\vfill
\newpage

% ==========================================================
% SLIDE 9: THE BIG TAKEAWAY
% ==========================================================

\vspace*{0.35in}

\begin{center}

{\fontsize{14}{16}\selectfont\color{brandblue}\textbf{THE ENGINEERING TAKEAWAY}}

\vspace{0.3in}

\begin{minipage}{0.88\textwidth}
\centering

{\color{brandblue}\rule{6cm}{1pt}}

\vspace{0.2in}

{\fontsize{24}{28}\selectfont\bfseries\color{brandnavy}
Sometimes the best algorithm\\
isn't the one that is theoretically\\
optimal.
}

\vspace{0.15in}

{\fontsize{12}{14}\selectfont\color{brandgray}
It is the one that best exploits the underlying hardware.

\vspace{0.1in}
The biggest leaps in AI didn't come from forcing networks to read better. They came from restructuring the math so GPUs could compute it faster.
}

\vspace{0.2in}

{\color{brandblue}\rule{6cm}{1pt}}

\end{minipage}

\end{center}

\vfill

\begin{center}

\rule{0.6\textwidth}{0.5pt}

\vspace{0.12in}

{\fontsize{11}{13}\selectfont\color{brandgray}
Building AI systems by understanding the mathematics,\\
engineering, and intuition behind every architecture.
}

\vspace{0.12in}

{\fontsize{10}{12}\selectfont\color{brandgray}
\begin{tabular}{@{}l l}
GitHub: & \texttt{github.com/MrHeaven1y} \\
Portfolio: & \texttt{mrheaven1y.github.io} \\
LinkedIn: & \texttt{linkedin.com/in/dibayendu-mukherjee-bb897b267}
\end{tabular}
}

\vspace{0.12in}

{\fontsize{14}{16}\selectfont\color{brandnavy}\textbf{Dibayendu Mukherjee}}

\end{center}

\vspace*{0.2in}

\end{document}
